% META: {"id":"1NSI-TC-CO-004","chapitre":"1NSI-TYPES-CONSTRUITS","type_objet":"corrige","capacites":["1NSI-TYPES-CONSTRUITS-C1"],"mode_creation":"ex_nihilo","sources_inspiration":[],"status":"approved","exercice_ref":"1NSI-TC-EX-004","origin":"needs_review","status_history":["needs_review","audited","accepted","approved"],"release_acceptance":"RELEASE_OWNER_BATCH_ACCEPTANCE","acceptance_closure_digest":"sha256:9b3ccf9a81c5520fbb7e03b7b2d3e7bf2057a02834b6d908bf3be5f49f3b553f","acceptance_receipt_digest":"sha256:4fad86f0282e0fa4e0419cca1ad6379ae7af5a280c04a4a90880f90a1c044242"}
\begin{corrige}{1NSI-TC-EX-004}
\begin{enumerate}
  \item Les trois affichages sont : \lstinline{Fatou}, puis \lstinline{True} (car $2350 > 2000$), puis \lstinline{diamant} (élément d'indice 2).

  \item C'est le \textbf{déballage de tuple} (\textit{tuple unpacking}) : Python affecte chaque élément du tuple à la variable correspondante, dans l'ordre.

  \item Non. Un tuple est \textbf{non mutable} : l'instruction \lstinline{joueur[1] = 2400} provoque une erreur \lstinline{TypeError: 'tuple' object does not support item assignment}.
\end{enumerate}

Vérification : les résultats ont été confirmés par exécution.
\end{corrige}
